%!TEX TS-program = lualatex
%!TEX encoding = UTF-8 Unicode

\documentclass[12pt, addpoints]{exam}

%\printanswers

\usepackage[left=1in,right=0.75in,top=1in]{geometry}     

\usepackage{fontspec}
\def\mainfont{Linux Libertine O}
\setmainfont[Ligatures={TeX, Common}, BoldFont={* Bold}, ItalicFont={* Italic}, Fractions ={On}, Numbers={OldStyle}]{\mainfont}
%\setmonofont[Scale=MatchLowercase]{Inconsolata} 
\setsansfont[Scale=MatchLowercase]{Linux Biolinum O} 
\usepackage{microtype}

\usepackage{graphicx}
	\graphicspath{%
	{/Users/goby/Pictures/teach/434/exams/}}%

% Used to get the last two digits of the year for the header below.
\def\short#1{\csname @gobbletwo\expandafter\endcsname\number#1}

\pagestyle{headandfoot}
\firstpageheader{Marine Evolutionary Ecology, Exam 2, \textsc{f}\short{\year}}{}{%
	\ifprintanswers\textbf{KEY, \numpoints\,pts.}\else Name: \enspace \makebox[2.5in]{\hrulefill}\fi}
\runningheader{}{}{\small{pg. \thepage}}
%\runningheadrule

\footer{}{}{\iflastpage{\small \rule{0.6in}{0.4pt} / \numpoints\ points}{}}%{\makebox[0.75in]{\hrulefill}}
\extrafootheight{-0.5in}

\pointsinmargin
\marginbonuspointname{ pts \textsc{e.c.}}
\marginpointname{ pts}


\usepackage{multicol}

%% Remove comment %% to print answer key.
\unframedsolutions
\renewcommand{\solutiontitle}{}
\SolutionEmphasis{\bfseries}

%% Create a Matching question format
\newcommand*\Matching[1]{
\ifprintanswers
	\textbf{#1}
\else
	\rule{2in}{0.5pt}
\fi
}
\newlength\matchlena
\newlength\matchlenb
\settowidth\matchlena{\rule{2.1in}{0pt}}
\newcommand\MatchQuestion[2]{%
	\setlength\matchlenb{0.98\linewidth}
	\addtolength\matchlenb{-\matchlena}
	\parbox[t]{\matchlena}{\Matching{#1}}\enspace\parbox[t]{\matchlenb}{#2}}


\newcommand*\AnswerBox[2]{%
    \parbox[t][#1]{0.92\textwidth}{%
    \begin{solution}#2\end{solution}}
    \vspace*{\stretch{1}}
}

\newenvironment{AnswerPage}[1]
    {\begin{minipage}[t][#1]{0.92\textwidth}%
    \begin{solution}}
    {\end{solution}\end{minipage}
    \vspace*{\stretch{1}}}

\newlength{\basespace}
\setlength{\basespace}{5\baselineskip}


\begin{document}

\noindent I score your exam based on what you write. Show me that you know the material. Use appropriate vocabulary. Rock on! Rock lobster!

\begin{questions}

%
\question[5]
Is larval duration (e.g., 14 days or 60 days) a good measure of how far larvae \emph{actually} disperse? Why?

	\AnswerBox{0.6\basespace}{No. Although larvae may have the potential to disperse far, they may not actually do so. Several studies have shown that some species with high potential do not disperse.}

%
\question
\begin{parts}
\part[5] Are two populations genetically connected or isolated if \textit{F}\textsubscript{ST} = 0.46? Are they ecologically connected or isolated? Explain.

	\AnswerBox{0.6\basespace}{This is a moderately high \textit{F}\textsubscript{ST} indicating little larval dispersal or gene flow. Therefore, the populations should be ecologically (2 pts) and genetically isolated (2 pts).}
	
	\part[5] Is each of these two populations more likely to be sustained by larval retention or immigration? Explain.
	
	\AnswerBox{0.6\basespace}{If there is little larval connectivity, then each population is most likely sustained through larval retention.}
\end{parts}

%
\question[5]
\begin{multicols}{2}
Why does the kelp forest distribution extend farther north on the west coast of Africa (dark shading, right) than on the east coast? 

\ifprintanswers\textbf{Kelp forests require cold water. The current gyres in the southern hemisphere carry cold water farther north on the west coast than on the east coast, allowing the kelp forest to extend farther north.}\fi

\columnbreak

\hfil\includegraphics[width=0.3\textwidth]{exam2_africa}\hfill
\end{multicols}


\newpage

%
\question[10]
Why is rocky intertidal ecosystem function highest with intermittent upwelling instead of constant upwelling or downwelling?

%% Change the question to ask why top down effects of predators are not important under states of constant upwelling or downwelling. 
% E.g., Why is rocky intertidal ecossytem function higher with intermittent upwelling instead of constant upwelling? Explain this from a bottom-up and top down perspective. (LAST part not great but a start).

	\begin{AnswerPage}{0.4\textheight}%
	 Relaxation of upwelling allows recruitment into the intertidal. (3 points)
	 
	 \vspace*{\baselineskip}
	 
	 With constant upwelling, water moves away from shore, reducing recruitment. (3 points)
	 
	 \vspace*{\baselineskip}
	 
	 With constant downwelling, larvae tend to be carried below the surface, reducing recruitment. (3 points)
	
	\end{AnswerPage}

%
\question[10]
Why does an alternate stable state (e.g., state \textsc{b}) not return to its originate state (e.g., state \textsc{a}) when environmental conditions return to their original state? Include an explanation of hysteresis as part of your answer. You may use a diagram to aid your answer.


	\begin{AnswerPage}{0.4\textheight}%
	 Hysteresis is the difference in conditions required to return shift between alternate stable states. (4 points).
	 
	\vspace*{\baselineskip}
		
	Thus, once one state flips to another state, it takes a very dramatic change in environmental conditions to flip the state back to its original state. (4 points)
	
	\end{AnswerPage}

\newpage

%
\question[5]
What are trait-mediated indirect interactions (\textsc{tmii})? Explain it but do not give me an example.
% Tie in Inducible defenses.

\AnswerBox{0.2\textheight}{\textsc{Tmii}s are the indirect effects that a top predator has on abundance of lower trophic level organisms by through the modification of prey's behavioral or morphological traits due to the presence of the predator. }

%

\question[10]
Explain whether the amount of trophic heat in the ecosystem would increase or decrease if climate change causes the greatest temperature difference between immersed and exposed \textit{Pisaster} sea stars.

\begin{AnswerPage}{0.6\textheight}%
	A study showed that feeding rates decreased the most in \textit{Pisaster} when sea stars experienced the greatest thermal stress (greatest temp difference). If feeding rates go down, then trophic heat will increase because they energy they have consumed is lost as metabolic heat instead of being passed to higher trophc levels.	
\end{AnswerPage}

\newpage

%

\bonusquestion[2]
For one extra credit point, what plankton group is 0.02–0.2 mm in size? For one more extra credit point, write the name of your favorite movie.

\AnswerBox{\baselineskip}{Microplankton (1 pt) / Whatever (1 pt)}

\question[10]
\begin{multicols}{2}
	The relative abundance of micro-, nano-, and picophytoplankton varies along a gradient from cold temperate to tropical latitudes. Explain how the figure at right provides support the change in relative abundance along the gradient, using your knowledge of competitive abilities among the phytoplankton groups and nutrient availability and high and low latitudes.
	
	\columnbreak
	
	\includegraphics[width=0.49\textwidth]{exam2_picoplankton}
	
\end{multicols}	

\begin{AnswerPage}{0.5\textheight}
	
	Tropical latitudes are nutrient poor so support few phytoplankton. Picoplankton are more efficient competitors in warm, nutrient poor waters. The low total Chl~\textit{a} (x-axis) corresponds to tropical waters, where picoplankton have the higher relative abundance. High total Chl~\textit{a} corresponds to colder temperate water where microphytoplankton (e.g., diatoms) have higher relative abundance.
	
\end{AnswerPage}

%

\newpage

\question
Whalen et al.~(2013) tested the mutualistic mesograzers hypothesis (\textsc{mmh}) in Virginia seagrass beds. Treatment plots (solid circles) received a chemical deterrent for Isopod mesograzers. Control plots (open circles) did not receive the deterrent. Mesograzer density (left panel) and amount of epiphytes (chlorophyll \textit{a}; right panel) were recorded before adding the deterrent (asterisk). The density of mesograzers and amount of epiphytes were recorded over several days.

\includegraphics[width=0.96\textwidth]{exam2_mesograzers}
\begin{parts}
	\part[5] Explain the mutualistic mesograzers hypothesis and how it affects the abundance of seagrasses.
	
	\AnswerBox{\basespace}{Mutalistic mesograzers hypothesis states that mesograzers consume epiphytes that grow on seagrasses, decreasing their functional efficiency. The mesograzers provide an indirect effect on seagrass abundance.}
	
	\part[5] Explain whether the results in the figure above provide support for or falsify the \textsc{mmh} and why. 

	\AnswerBox{\basespace}{The figure above does provide support. As mesograzers decline in density, the amount of epiphytes increases.}

\end{parts}



\end{questions}


\end{document}